% vitae-doc.tex -- user manual of the vitae class (built with `l3build doc`).
% Samples (sample-*.tex) are complete vitae documents, compiled first and
% shown here as source beside their output.
\documentclass{l3doc}
\ProvidesFile{vitae-doc.tex}[2026-11-22 v0.4.0 vitae user manual]
\usepackage{fancyvrb,graphicx,booktabs,array}

% \sample{<name>}: the source of <name>.tex, then its cropped output.
\newcommand\sample[1]{%
  \par\medskip
  % gobble=0: l3doc (via doc) would strip two characters from each line.
  \VerbatimInput[gobble=0, frame=single, fontsize=\small, framesep=2mm,
    label={\normalfont\small\ttfamily #1.tex}]{#1.tex}%
  \par\nopagebreak\smallskip
  \begin{center}
    \fbox{\includegraphics[width=0.92\linewidth]{#1.pdf}}
  \end{center}\medskip}
\providecommand\XeLaTeX{Xe\LaTeX}
% \key{<name>}: a key in running text.
\newcommand\key[1]{\texttt{#1}}

\GetFileInfo{vitae-doc.tex}
\title{The \pkg{vitae} class\thanks{This manual describes version \fileversion, dated \filedate.}}
\author{\url{https://github.com/gvgramazio/tex-resume}}
\date{Released \filedate}

\begin{document}
\maketitle

\begin{abstract}
\noindent
\pkg{vitae} is a \LaTeX{} class for resumes. You declare your data once, in
the preamble or in a separate file, and a \emph{theme} decides how it looks,
in the same way \pkg{beamer} separates content from presentation. It works
with pdf\LaTeX, \XeLaTeX{} and Lua\LaTeX, produces text that applicant-tracking
systems extract cleanly, and can produce accessible (PDF/UA) and archival
(PDF/A) files.
\end{abstract}

\tableofcontents

\section{Quick start}

A resume is a list of declarations followed by a body that prints them.
\cs{makeheader} prints the name and contacts, as \cs{maketitle} does for a
title, and \cs{printentries} prints a list of entries under a heading.

\sample{sample-quickstart}

Every vitae document follows this pattern. The data can live in its own file
(\cs{input}\{data\}) so that several documents, in different themes or
languages, share it.

\section{Personal data and contacts}

\begin{function}{\name,\firstname,\lastname,\jobtitle,\photo}
  \begin{syntax}
    \cs{name}\marg{first}\marg{last}
    \cs{jobtitle}\marg{text}
    \cs{photo}\marg{file}
  \end{syntax}
  Identity. \cs{name} sets \cs{firstname} and \cs{lastname} at once. The photo
  is a graphics file, with or without extension; a missing file only warns.
  Its width is the \key{photowidth} setting.
\end{function}

\begin{function}{\email,\phone,\address,\homepage,\linkedin,\github,\orcid}
  \begin{syntax}
    \cs{email}\marg{address} \quad \cs{homepage}\marg{example.com/path}
    \cs{linkedin}\marg{id} \quad \cs{github}\marg{user} \quad \cs{orcid}\marg{id}
  \end{syntax}
  Contacts, printed by \cs{makeheader} in this fixed order (email, phone,
  address, website, LinkedIn, GitHub, ORCID), each with an icon and, except
  for phone and address, a link. Values that are addresses or handles may
  contain |_ & ~ #| as typed, and also accept |\_ \& \# \%|; a raw |%| starts
  a comment, as everywhere in \LaTeX.
\end{function}

\begin{function}{\addcontact}
  \begin{syntax}
    \cs{addcontact}\oarg{url}\marg{icon}\marg{text}
    \cs{addcontact}\oarg{options}\marg{icon}\marg{text}
  \end{syntax}
  Any other contact, printed after the built-in ones in the order given.
  The options are \key{url} and \key{label}; the \key{label} (e.g.\ |Blog|) is what text extraction and screen readers
  read for the icon; without it the icon is simply hidden from them. Brace
  values that contain |=| or |,|, as in |url={https://example.org/?a=1}|.
\end{function}

\begin{function}{\thefirstname,\ifvitaeset,\makeheader,\printcontacts,\printphoto}
  \begin{syntax}
    \cs{the}\meta{field} \quad \cs{ifvitaeset}\marg{field}\marg{true}\marg{false}
  \end{syntax}
  Every field \meta{field} above has an expandable getter, such as
  \cs{theemail}, and \cs{ifvitaeset} tests whether it is set. \cs{makeheader}
  prints the header template; \cs{printcontacts} and \cs{printphoto} are its
  building blocks, for use in templates.
\end{function}

\section{Entries, items and lists}

\begin{function}{\cventry}
  \begin{syntax}
    \cs{cventry}\marg{keys}\marg{body}
  \end{syntax}
  An entry: a job, a degree, an award. All keys are optional:
  \begin{description}
    \item[\key{title}, \key{org}, \key{location}] the entry's text fields;
    \item[\key{start}, \key{end}] its dates (see section~\ref{sec:dates});
      an empty \key{end} means ongoing;
    \item[\key{date}] a single date: sets \key{start} and \key{end} at once;
    \item[\key{style}] the layout of this entry (section~\ref{sec:styles});
    \item[\key{dateformat}, \key{datelocale}, \key{durations}] date settings
      for this entry only.
  \end{description}
  The body is free text; blank lines separate paragraphs, spaced by the
  \key{parskip} setting.
\end{function}

\begin{function}{\cvitem}
  \begin{syntax}
    \cs{cvitem}\marg{label}\marg{text}
  \end{syntax}
  A one-line item, for skills or languages.
\end{function}

\begin{function}{\addentry,\additem,\printentries}
  \begin{syntax}
    \cs{addentry}\marg{list}\marg{keys}\marg{body}
    \cs{additem}\marg{list}\marg{label}\marg{text}
    \cs{printentries}\oarg{heading}\marg{list}
    \cs{printentries}\oarg{options}\marg{list}
  \end{syntax}
  Named lists store entries and items in the order they were added, so data
  and layout stay apart: the same list can be printed several times, in
  different documents or styles. \cs{printentries} prints a list under a
  section heading, but only if the list has content; an empty or undefined
  list warns and prints nothing. Its optional argument is either a plain
  heading or options: \key{heading}, \key{style} and \key{durations}. It
  counts as options if it contains |=| or starts with an option name, as in
  |[durations]|; brace a heading that contains a comma, as in
  |heading={Talks, Panels}|.
\end{function}

\section{Entry styles}\label{sec:styles}

Each theme offers three layouts for entries. \key{full} is the default;
\key{compact} puts an entry on one line, for certifications or awards;
\key{inline} is a run-in paragraph, for talks or publications. Choose one for
a whole list with \cs{printentries}|[style=...]|, or for one entry with the
\key{style} key, which wins over the list's.

\sample{sample-styles}

A style \meta{name} is the template |entry/|\meta{name}; \cs{settemplate}
overrides one, or defines a new style (section~\ref{sec:templates}). An
unknown style warns once and falls back to \key{full}.

\section{Dates}\label{sec:dates}

Dates are written as ISO dates, |2020|, |2020-03| or |2020-03-01|, or as any
other text, such as |Spring 2016|. ISO dates are printed in the current
\emph{style} and \emph{locale}; any other text is printed exactly as typed.

\begin{function}{\formatdate,\daterange,\duration}
  \begin{syntax}
    \cs{formatdate}\marg{date} \quad \cs{daterange}\marg{start}\marg{end}
    \cs{duration}\marg{start}\marg{end}
  \end{syntax}
  All three are expandable. \cs{daterange} prints \enquote{start -- end};
  without an end, \enquote{start -- present}, with the locale's word unless
  the \key{present} setting gives one; equal dates print once.
  \cs{duration} prints the length of a range, such as \enquote{2 yrs 7 mos}.
\end{function}

\subsection{Styles and locales}

The \key{dateformat} setting picks the style: \key{long} (the default:
\enquote{March 2020}, \enquote{March 1, 2020}), \key{short} (\enquote{Mar
2020}), \key{numeric} (\enquote{03/2020}), \key{iso} (\enquote{2020-03}) or
\key{year} (\enquote{2020}). A year on its own always prints as the year.

The locale follows the document language, from \pkg{babel} or
\pkg{polyglossia}, including \cs{selectlanguage}; the \key{datelocale}
setting forces one. Built-in locales are \key{english} (US order),
\key{british}, \key{italian}, \key{german}, \key{french} and \key{spanish},
with the usual \pkg{babel} aliases (\key{ngerman}, \key{UKenglish},
\ldots); any other language falls back to English with a warning.

\begin{function}{\vitaedatelocale}
  \begin{syntax}
    \cs{vitaedatelocale}\marg{name}\marg{keys}
  \end{syntax}
  Defines a locale, starting from English for every key not given:
  \key{months} and \key{short} (twelve names each), \key{present}, \key{units}
  (|{yr, yrs, mo, mos}|: year and month, singular and plural), \key{aliases}
  (other \cs{languagename} values), and the patterns \key{ym}, \key{ymd},
  \key{numeric-ym} and \key{numeric-ymd}, which receive |#1| = year, |#2| =
  month and |#3| = day.
\end{function}

\noindent For example:
\begin{Verbatim}[gobble=0, fontsize=\small]
\vitaedatelocale{portuguese}{
  months  = {janeiro, fevereiro, março, abril, maio, junho, julho,
             agosto, setembro, outubro, novembro, dezembro},
  present = atual, units = {ano, anos, mês, meses},
  ym = {#2 de #1}, ymd = {#3 de #2 de #1}}
\end{Verbatim}

\subsection{Durations}

With the \key{durations} setting, each entry also shows the length of its
range, next to its dates; it can be enabled for the whole document, for one
list (\cs{printentries}|[durations]|) or for one entry (the \key{durations}
key). Months count inclusively, so September 2018 -- March 2021 is
\enquote{2 yrs 7 mos}. If either date is a year only, the length is in
whole years (2018 -- 2021 is \enquote{3 yrs}) and a range within one year
has none. Single dates, free text and an end before the start give none.
An ongoing range ends \emph{today}: the date of each compilation, or the
\key{today} setting, which fixes it for reproducible output.

\sample{sample-dates}

\subsection{Scope}

Date settings, like every setting in this manual, are local: in the preamble
they apply to the whole document, inside a group or environment only there,
and as entry keys to one entry.

\section{Themes and settings}

\begin{function}{\usetheme}
  \begin{syntax}
    \cs{usetheme}\marg{name}
  \end{syntax}
  Loads the theme |vitaetheme|\meta{name}|.sty|, once, in the preamble.
  Without \cs{usetheme}, \key{classic} is used. Two themes ship with vitae:
  \key{classic} (serif, centred header, dates on the right, as in the
  samples above) and \key{modern} (sans serif, accent colour, dates in a
  margin column, photo).
\end{function}

\sample{sample-modern}

\begin{function}{\vitaesetup}
  \begin{syntax}
    \cs{vitaesetup}\marg{keys}
  \end{syntax}
  General settings. User settings always win over theme defaults, whether
  they are made before or after \cs{usetheme}.
\end{function}

\begin{center}
\begin{tabular}{@{}l>{\raggedright\arraybackslash}p{0.62\linewidth}@{}}
  \toprule
  Key & Meaning (default) \\
  \midrule
  \key{accent} & the theme's accent colour, any \pkg{xcolor} colour \\
  \key{parskip} & space between paragraphs of free text and of entry
    bodies (|0.5\baselineskip plus 2pt|; |0pt| turns it off) \\
  \key{photowidth} & width of the photo (|2.5cm|) \\
  \key{dateformat} & date style: \key{long}, \key{short}, \key{numeric},
    \key{iso}, \key{year} (\key{long}) \\
  \key{datelocale} & date locale (the document language) \\
  \key{datesep} & separator in date ranges (a thin-spaced en dash) \\
  \key{present} & word for an ongoing range (the locale's) \\
  \key{durations} & show the duration of entries (off) \\
  \key{today} & the end of ongoing durations, an ISO date (compilation date) \\
  \bottomrule
\end{tabular}
\end{center}

\section{Templates}\label{sec:templates}

A theme is a set of \emph{templates}: |header|, |contact|, |contactsep|,
|section|, |entry|, the entry styles |entry/|\meta{style}, and |item|. Any of
them can be replaced without writing a theme.

\begin{function}{\settemplate,\field,\iffield,\ifanyfield,\joinfields,\addpunct}
  \begin{syntax}
    \cs{settemplate}\marg{name}\marg{code}
    \cs{field}\marg{field} \quad \cs{iffield}\marg{field}\marg{true}\marg{false}
    \cs{ifanyfield}\marg{field, \ldots}\marg{true}\marg{false}
    \cs{joinfields}\marg{separator}\marg{fields}
    \cs{addpunct}\marg{punctuation}
  \end{syntax}
  Inside a template, \cs{field} reads the current item's fields: |title|,
  |org|, |location|, |start|, |end|, |dates| (the formatted range),
  |duration|, |body| and |label|; and for contacts |icon|, |text| and |url|.
  \cs{joinfields} prints only the fields that are set, from a list such as
  |title=\textbf, org, location| (each with an optional format), so that missing fields leave no stray separator.
  \cs{addpunct} adds punctuation unless the text already ends with some.
  Defining |entry/|\meta{name} creates a new entry style.
\end{function}

\sample{sample-template}

Templates should not handle paragraph spacing: vitae runs them with
|\parskip| = 0 and spaces body paragraphs itself. For accessible output, a
template wraps headings, decoration and meaningful symbols in the helpers of
the next section.

\section{Accessible PDF and applicant-tracking systems}

Every vitae PDF extracts as clean text, which is what applicant-tracking
systems (ATS) read: contact icons come out as their label (\enquote{Email:
jane.doe@example.com}, \enquote{GitHub: jdoe-example}) and contacts never run
together. That needs no setup.

For a tagged, accessible PDF that conforms to PDF/UA, make this the first
line and compile with \textbf{Lua\LaTeX}:
\begin{Verbatim}[gobble=0, fontsize=\small]
\DocumentMetadata{tagging=on, pdfstandard=ua-2, lang=en-US}  % PDF/UA-2
% or: tagging=on, pdfstandard=ua-1, pdfversion=1.7, lang=en-US  % PDF/UA-1
\documentclass{vitae}
\end{Verbatim}
The name and section titles become headings, icons get replacement text,
decoration is marked as such, the photo gets alternative text, and the
document title is set from \cs{name} and \cs{jobtitle}.

\begin{function}{\taggedheading,\decorative,\actualtext}
  \begin{syntax}
    \cs{taggedheading}\marg{level}\marg{material}
    \cs{decorative}\marg{material}
    \cs{actualtext}\marg{text}\marg{content}
  \end{syntax}
  For templates. \cs{taggedheading} tags whole paragraphs as a heading of
  the given level; \cs{decorative} marks rules and other drawn decoration;
  \cs{actualtext} makes \meta{content} read, and extract, as \meta{text}
  (or as nothing, if empty). Without tagging, they only typeset their
  content.
\end{function}

\section{Archival PDF (PDF/A)}

PDF/A is the ISO standard for long-term archiving, which some employers, job
portals and public bodies ask for. Choose a recipe and make it the first
line of the document:

\begin{description}
  \item[Archival and accessible, PDF 2.0] Lua\LaTeX:\\
    |\DocumentMetadata{tagging=on, pdfstandard={a-4f,ua-2}}|
  \item[Archival and accessible, PDF 1.7] Lua\LaTeX; the widest support:\\
    |\DocumentMetadata{tagging=on, tagging-setup={attach-css=false},|\\
    |  pdfstandard={a-2a,ua-1}, pdfversion=1.7}|
  \item[Same, embedded files allowed] Lua\LaTeX:\\
    |\DocumentMetadata{tagging=on,|\\
    |  pdfstandard={a-3a,ua-1}, pdfversion=1.7}|
  \item[Archival, all text as Unicode] pdf\LaTeX{} or Lua\LaTeX:\\
    |\DocumentMetadata{pdfstandard=a-2u, pdfversion=1.7}|
  \item[Archival, visual fidelity only] any engine:\\
    |\DocumentMetadata{pdfstandard=a-2b, pdfversion=1.7}|
\end{description}

Add |lang=| with the document language to each. Known limitations:
\begin{itemize}
  \item PDF/A-2 and PDF/A-3 need |pdfversion=1.7|: \LaTeX{} writes PDF 2.0
    by default, which only PDF/A-4 accepts (as for PDF/UA-1).
  \item PDF/A-2 with tagging needs |attach-css=false|: \LaTeX's tagging
    embeds two HTML stylesheets, and PDF/A-2 forbids embedded files that are
    not PDF/A themselves. PDF/A-3 and PDF/A-4 allow them.
  \item The accessible (\enquote{a}) levels need tagging, and so Lua\LaTeX.
  \item \XeLaTeX{} cannot reach PDF/A-2u: its icon fonts have no Unicode map.
  \item PDF/A-1 is not tested. A photo should be a JPEG, a PNG or a PDF/A
    file: images are embedded as they are.
\end{itemize}

\section{Engines}

\begin{center}
\begin{tabular}{@{}lccc@{}}
  \toprule
  & pdf\LaTeX & \XeLaTeX & Lua\LaTeX \\
  \midrule
  Everything in this manual, untagged & yes & yes & yes \\
  Clean text extraction (ATS), with labels & yes & without labels & yes \\
  Tagged PDF (PDF/UA), PDF/A \enquote{a} levels & no (warns) & no (warns) & yes \\
  PDF/A-2u & yes & no & yes \\
  \bottomrule
\end{tabular}
\end{center}

\noindent On \XeLaTeX{} the icons are hidden from extracted text rather than
labelled, because its PDF driver would place the labels out of reading order.

\section{Installation}

\pkg{vitae} needs \TeX{} Live 2022 or later (or an equivalent MiK\TeX).
Download |vitae-ctan.zip| from the latest release on GitHub and copy
|vitae.cls| and the |vitaetheme*.sty| files either next to your document or
into your personal tree (|~/texmf/tex/latex/vitae| on Linux,
|~/Library/texmf/tex/latex/vitae| on macOS). From a clone of the repository,
|l3build install| does the same.

\end{document}
